\newlength{\PageHeight}
\newlength{\PageWidth}
% ==============================================
% 一、基础框架 & 页面布局
% ==============================================
%\documentclass[17pt]{beamer}
%\setlength{\PageHeight}{13.5cm}
%\setlength{\PageWidth}{24cm} 

\documentclass[20pt]{beamer}
\setlength{\PageHeight}{19.05cm}
\setlength{\PageWidth}{33.866cm} 

% 设置页面边距变量为长度寄存器
\newlength{\PageLeft}
\newlength{\PageRight}
\newlength{\PageTop}
\newlength{\PageBottom}

\setlength{\PageLeft}{0.10498\PageHeight}
\setlength{\PageRight}{0.10498\PageHeight} 
\setlength{\PageTop}{0.07874\PageHeight} 
\setlength{\PageBottom}{0.10498\PageHeight}

% 自定义宽屏尺寸
\usepackage{geometry}
\geometry{
    paperwidth=\PageWidth,
    paperheight=\PageHeight,
    left=\PageLeft, right=\PageRight,
    top=\PageTop, bottom=\PageBottom
}

% 关闭默认导航、侧边栏、页眉
\setbeamertemplate{navigation symbols}{}
\setbeamertemplate{headline}{}
\setbeamertemplate{sidebar left}{}
\setbeamertemplate{sidebar right}{}

% 全局背景
\pagecolor{white}

% ==============================================
% 二、宏包加载（按依赖顺序）
% ==============================================
% 中日英字体 & 数学字体（XeLaTeX 专用顺序）
\usepackage{fontspec}
\usepackage{xeCJK}
\xeCJKsetup{CJKmath=true} % 支持公式内中文
\usepackage{unicode-math}

% 图形、颜色、绘图
\usepackage{graphicx}
\usepackage{xcolor}
\usepackage{tikz}
\usetikzlibrary{positioning}
\usetikzlibrary{backgrounds}
\usepackage{float}

% 表格
\usepackage{booktabs, multirow, bigstrut}

% 数学环境
\usepackage{amsfonts, amsmath, amsthm}
\usepackage{amssymb,bm,mathrsfs} % 数学符号支持


% 算法
\usepackage{algorithm, algpseudocode}

% 代码高亮
\usepackage{listings}

% 图表子图 & 标题
\usepackage{caption, subcaption}

% 多栏排版
\usepackage{multicol}

% ==============================================
% 三、全局段落格式
% ==============================================
\linespread{1.3}                 % 行距
\setlength{\parindent}{2em}      % 首行缩进
\BeforeBeginEnvironment{minipage}{\noindent}
\setlength{\parskip}{0.5\baselineskip}  % 段后0.5行
% ==============================================
% 四、字体配置（思源系列 + 英文字体）
% ==============================================
% 中文无衬线字体（思源黑体）
\setCJKsansfont[
    Path              = ./font/,
    Extension         = .otf,
    UprightFont       = SourceHanSansCN-Regular,
    BoldFont          = SourceHanSansCN-Bold,
    ItalicFont        = SourceHanSansCN-Regular,
    BoldItalicFont    = SourceHanSansCN-Bold,
    AutoFakeBold      = false,
    ItalicFeatures    = {FakeSlant = 0.15},
    BoldItalicFeatures= {FakeSlant = 0.15},
    Mapping           = fullwidth-stop% 把中文句号自动转化英文小圆点
]{}
% \setCJKsansfont{Microsoft YaHei} % 也可替换为微软雅黑（但需注意版权问题）

% 中文衬线字体（思源宋体）
\setCJKmainfont[
    Path              = ./font/,
    Extension         = .otf,
    UprightFont       = SourceHanSerifCN-Regular,
    BoldFont          = SourceHanSerifCN-Bold,
    ItalicFont        = SourceHanSerifCN-Regular,
    BoldItalicFont    = SourceHanSerifCN-Bold,
    AutoFakeBold      = false,
    ItalicFeatures    = {FakeSlant = 0.15},
    BoldItalicFeatures= {FakeSlant = 0.15},
    Mapping           = fullwidth-stop% 把中文句号自动转化英文小圆点
]{}
%\setCJKmainfont{SimSun} % 也可替换为宋体（但需注意版权问题）

% 中文等宽字体（霞鹜文楷）
\setCJKmonofont[
    Path              = ./font/,
    Extension         = .ttf,
    UprightFont       = LXGWWenKaiMono-Regular,
    BoldFont          = LXGWWenKaiMono-Medium,
    ItalicFont        = LXGWWenKaiMono-Regular,
    BoldItalicFont    = LXGWWenKaiMono-Medium,
    AutoFakeBold      = false,
    ItalicFeatures    = {FakeSlant = 0.15},
    BoldItalicFeatures= {FakeSlant = 0.15},
    Mapping           = fullwidth-stop% 把中文句号自动转化英文小圆点
]{}
%\setCJKmonofont{KaiTi} % 也可替换为楷体（但需注意版权问题）

% 英文字体
\setmainfont{TeX Gyre Termes}
\setsansfont{TeX Gyre Heros}
\setmathfont{TeX Gyre Termes Math} % 或选择 Latin Modern Math
\setmonofont[
    Path              = ./font/,
    Extension         = .otf,
    UprightFont       = SourceCodePro-Regular,
    BoldFont          = SourceCodePro-Bold,
    ItalicFont        = SourceCodePro-It,
    BoldItalicFont    = SourceCodePro-BoldIt,
    AutoFakeBold      = false
]{}
%\setmonofont{TeX Gyre Cursor} %等宽也可以替换成这个字体


% ==============================================
% 五、全局色彩定义
% ==============================================
\definecolor{NPUBlue}{RGB}{0,78,162}
\definecolor{NPUDarkBlue}{RGB}{0,51,102}
\definecolor{ExampleGreen}{RGB}{0,128,64}
\definecolor{AlertOrange}{RGB}{204,102,0}
\definecolor{LightTea}{RGB}{245,241,236}

% ==============================================
% 六、Beamer 全局颜色样式
% ==============================================
% 列表符号颜色
\setbeamercolor{item}{fg=NPUBlue}
\setbeamercolor{subitem}{fg=NPUBlue}
\setbeamercolor{subsubitem}{fg=NPUBlue}
\setbeamercolor{enumerate item}{fg=NPUBlue}

% 各级标题颜色
\setbeamercolor{frametitle}{fg=NPUBlue}
\setbeamercolor{framesubtitle}{fg=NPUBlue}
\setbeamercolor{section title}{fg=NPUBlue}
\setbeamercolor{subsection title}{fg=NPUBlue}

% 线条、脚注、标签
\setbeamercolor{separation line}{fg=NPUBlue}
\setbeamercolor{footnote}{fg=NPUBlue}
\setbeamercolor{label}{fg=NPUBlue}

% 图表标题汉化 + 颜色
\renewcommand{\figurename}{图：}
\renewcommand{\tablename}{表：}
\captionsetup{labelfont = {color=NPUBlue}, labelsep=space}

% 参考文献样式与颜色
\setbeamertemplate{bibliography item}[text]
\setbeamercolor{bibliography item}{fg=NPUBlue}
\setbeamercolor{bibliography entry author}{fg=NPUBlue}
\setbeamercolor{bibliography entry title}{fg=black}
\setbeamercolor{bibliography entry location}{fg=NPUBlue!60}

% ==============================================
% 七、区块环境（块、定理、定义、示例、警示等）
% ==============================================
% 基础区块外观：圆角、无阴影、标题加粗
\setbeamertemplate{blocks}[rounded][shadow=false]
\setbeamerfont{block title}{series=\bfseries}

% 基础区块配色
\setbeamercolor{block title}{fg=white, bg=NPUBlue}
\setbeamercolor{block body}{fg=black, bg=NPUBlue!10}

% 示例区块（绿色）
\renewenvironment{exampleblock}[1]
{%
  \setbeamercolor{block title}{fg=white, bg=ExampleGreen}
  \setbeamercolor{block body}{fg=black, bg=ExampleGreen!10}
  \begin{block}{#1}
}
{\end{block}}

% 警示区块（橙色）
\renewenvironment{alertblock}[1]
{%
  \setbeamercolor{block title}{fg=white, bg=AlertOrange}
  \setbeamercolor{block body}{fg=black, bg=AlertOrange!10}
  \begin{block}{#1}
}
{\end{block}}

% 定理、引理、推论、命题
\renewenvironment{theorem}[1][]{\begin{block}{定理\ifx\relax#1\relax\else：#1\fi}}{\end{block}}
\renewenvironment{lemma}[1][]{\begin{block}{引理\ifx\relax#1\relax\else：#1\fi}}{\end{block}}
\renewenvironment{corollary}[1][]{\begin{block}{推论\ifx\relax#1\relax\else：#1\fi}}{\end{block}}
\newenvironment{proposition}[1][]{\begin{block}{命题\ifx\relax#1\relax\else：#1\fi}}{\end{block}}

% 定义（青色系）
\renewenvironment{definition}[1][]{
    \setbeamercolor{block title}{bg=cyan!70!black, fg=white}
    \setbeamercolor{block body}{bg=cyan!5}
    \begin{block}{定义\ifx\relax#1\relax\else：#1\fi}}{\end{block}}

% 例子
\renewenvironment{example}[1][]{\begin{exampleblock}{例子\ifx\relax#1\relax\else：#1\fi}}{\end{exampleblock}}

% 习题（紫色系）
\newenvironment{exercise}[1][]{
    \setbeamercolor{block title}{bg=purple!70!black, fg=white}
    \setbeamercolor{block body}{bg=purple!5}
    \begin{block}{习题\ifx\relax#1\relax\else：#1\fi}}{\end{block}}

% 注记（复用警示块）
\newenvironment{remark}[1][]{\begin{alertblock}{注记\ifx\relax#1\relax\else：#1\fi}}{\end{alertblock}}

% ==============================================
% 八、代码 & 算法样式
% ==============================================
\lstset{
    basicstyle=\ttfamily\small,
    columns=fixed,
    breaklines=true,
    frame=single,
    backgroundcolor=\color{gray!5},
    commentstyle=\color{gray},
    keywordstyle=\color{blue}\bfseries,
    stringstyle=\color{red},
    showstringspaces=false
}

\renewcommand{\algorithmname}{算法}
\algrenewcommand\algorithmicrequire{\textbf{输入:}}
\algrenewcommand\algorithmicensure{\textbf{输出:}}
\algrenewcommand\algorithmicreturn{\textbf{返回}}

% ==============================================
% 九、帧标题 自定义顶部栏
% ==============================================

\makeatletter
\setbeamertemplate{frametitle}{
  \begin{tikzpicture}[remember picture,overlay]
    \node[anchor=north west, inner sep=0pt]
        at ([yshift=-0.01994\PageHeight, xshift=0]current page.north west)
        {
\includegraphics[height=0.0944\PageHeight]{src/figure/frametitle_logo.jpg}\hspace{0.01049\PageHeight}};
    \node[anchor=north west, inner sep=0pt]
        at ([yshift=-0.04094\PageHeight, xshift=0.10498\PageHeight]current page.north west)
        {\fontsize{0.06641\PageHeight}{0.07379\PageHeight}\bfseries\color{NPUBlue}\insertframetitle};
  \end{tikzpicture}
}
\makeatother

% ==============================================
% 十、底部蓝色通栏页脚
% ==============================================
\setbeamertemplate{footline}{
    \newdimen\FootNoteMargin
    \FootNoteMargin=\PageTop
    \advance\FootNoteMargin by 0.05249\PageHeight
  \begin{overlayarea}{\paperwidth}{\the\FootNoteMargin}
    \usebeamertemplate***{footline rule}
    \usebeamertemplate***{footnote text}
  \end{overlayarea}

  \begin{tikzpicture}[remember picture, overlay]
    \begin{pgfonlayer}{background}
      \fill[NPUBlue]
        (current page.south west) rectangle
        ([yshift=0.03779\PageHeight]current page.south east);
    \end{pgfonlayer}
  \end{tikzpicture}

  \begin{tikzpicture}[remember picture, overlay]
    \tikzset{foottext/.style={font=\fontsize{0.02213\PageHeight}{0.02213\PageHeight}\selectfont\color{white}, inner sep=0pt}}
    \node[foottext, anchor=west]
      at ([xshift=0.10498\PageHeight, yshift=0.01889\PageHeight]current page.south west)
      {\ReportAuthor (\ReportInstitute)};
    \node[foottext, anchor=center]
      at ([yshift=0.01889\PageHeight]current page.south)
      {\ReportTitle};
    \node[foottext, anchor=east]
      at ([xshift=-0.10498\PageHeight, yshift=0.01889\PageHeight]current page.south east)
      {\insertframenumber\ / \inserttotalframenumber};
  \end{tikzpicture}
}



% ==============================================
% 十二、自定义页面命令：封面页、目录页
% ==============================================
% 封面标题页
\newcommand{\maketitlepage}{
    \begin{frame}[plain]
        \begin{tikzpicture}[remember picture,overlay]
            \node[anchor=north west, inner sep=0pt]
                at ([xshift=0.10498\PageHeight, yshift=-0.05498\PageHeight]current page.north west)
                {
\includegraphics[width=0.47244\PageHeight]{src/figure/logo_main.png}};
			\node[anchor=south, inner sep=0pt]
                at (current page.south)
                {
\includegraphics[width=\paperwidth]{src/figure/titlepage_cover.jpg}};
            \node[align=center, font=\Large, black]
                at (current page.center) {
                    {\fontsize{0.06642\PageHeight}{0.06642\PageHeight}\bfseries\color{NPUBlue} \ReportTitle     }\\[0.07874\PageHeight]
                    {\fontsize{0.04797\PageHeight}{0.04797\PageHeight}\bfseries                \ReportAuthor    }\\[0.0005\PageHeight]
                    {\fontsize{0.04797\PageHeight}{0.04797\PageHeight}\bfseries                \ReportInstitute }\\[0.02199\PageHeight]
                    {\fontsize{0.03321\PageHeight}{0.03321\PageHeight}\selectfont              \ReportTime      }
                };
        \end{tikzpicture}
    \end{frame}
}

% 目录样式：方形编号
\defbeamertemplate{section in toc}{customsquare}{%
  \raisebox{-0.2\height}{%
    \begin{tikzpicture}
      \node[
        fill=NPUBlue,
        minimum size=1cm,
        inner sep=0pt,
        outer sep=0pt,
        align=center,
        text=LightTea
      ] {\inserttocsectionnumber};
    \end{tikzpicture}%
  }
  \hspace{0.01574\PageHeight}
  \inserttocsection\par
}
\setbeamercolor{section in toc}{fg=NPUBlue}
\setbeamertemplate{section in toc}[customsquare]
\setcounter{tocdepth}{1}

% 目录页面
\newcommand{\mytocpage}{
    \begin{frame}[plain]
        \begin{tikzpicture}[remember picture,overlay]
            \node[anchor=south, inner sep=0pt]
                at (current page.south)
                {\includegraphics[width=\paperwidth]{src/figure/toc_blue.jpg}};
            \node[anchor=north west, inner sep=0pt]
                at ([yshift=-0.15748\PageHeight, xshift=0.15498\PageHeight]current page.north west)
                {\fontsize{0.06641\PageHeight}{0.07378\PageHeight}\bfseries\color{NPUBlue} 目录};
        \end{tikzpicture}

        \vspace{\dimexpr 1cm - \PageTop\relax}
        {\fontsize{0.04796\PageHeight}{0.04796\PageHeight}\bfseries
            \begin{multicols}{2}
                \noindent\tableofcontents
            \end{multicols}
        }
    \end{frame}
}

% ==============================================
% 十三、章节自动过渡页
% ==============================================
\AtBeginSection[]{
\begingroup
\begin{frame}[noframenumbering,plain]
\begin{tikzpicture}[remember picture,overlay]
    \node[anchor=south, inner sep=0pt]
        at (current page.south)
        {\includegraphics[width=\paperwidth]{src/figure/toc_blue.jpg}};
\end{tikzpicture}

\begin{columns}[c]
% 左侧超大淡数字
\begin{column}{0.25\textwidth}
\centering
\begin{tikzpicture}
\node[
    text=NPUBlue!12, % 淡透明蓝色数字
    font=\bfseries,
    scale=9
]{\thesection};
\end{tikzpicture}
\end{column}

% 右侧标题与标语
\begin{column}{0.55\textwidth}
\raggedleft
% 章节大标题
{\bfseries\LARGE\color{NPUDarkBlue} {\fontsize{38pt}{38pt} \insertsection }} \par
\vspace{3mm}
% 装饰横线
{\color{NPUDarkBlue!70}\rule{0.8\linewidth}{2pt}}\par
\vspace{3mm}
% 底部标语
{\ttfamily\small\color{black!40} 隐姓埋名，为国铸剑}
\end{column}
\end{columns}
\vfill
\end{frame}
\endgroup
}